\chapter{Base Geometry and Multi-Point Designs}
\label{ch:base-geometry}

At the fixed base point of the main manuscript every agent's frame is a point of one principal
fiber, the independent links are the only transports, and the design is \(\{c_*\}\). This chapter
asks what additional structure appears when contexts vary, and which of it is forced by the one
principal bundle rather than declared. Two local sections of that bundle determine a unique
relative frame, and the active gauge image at a context remains one copy of \(G\) inside the passive
product of independent belief and model re-expressions
(\Cref{prop:geo-relative-frame,prop:geo-two-atlas-active-image}). Chosen connections define
covariant defects of the cross maps, whose vanishing is exactly parallelness, and bundle
triviality, graph flatness, connection flatness, and trivial loop holonomy are four logically
independent conditions (\Cref{prop:geo-covariant-defect-parallelness,prop:geo-four-notions}). On a
finite design, constraining one link to be shared across design points leaves exactly the
context-independent rechoices multiplied by a central twist common to both endpoints
(\Cref{prop:gen-shared-link-stabilizer}). The connections and defects are used by the base-direction
pullbacks of \Cref{ch:pullback-geometry}, the designs and section responses by the section-space
histories of \Cref{ch:section-space-histories}, and the scale maps over base maps by the covers of
\Cref{ch:neighborhood-parents}.

The connections of \Cref{def:geo-connections} are chosen data, and no curvature or transport is
inferred from the agent frames; the matrix defect laws are not asserted for nonlinear law fibers;
graph links are not identified with base transports except under
\Cref{hyp:geo-graph-base-transport}; and the finite constructions of this chapter supply no
probability law on a continuum space of sections. \status{NOT-CLAIMED} The continuum theory is
\Cref{open:prob-continuum-theory}, the realization of independently specified links by base
transport and the recovery of a connection from records are
\Cref{open:rg-cover-graph-base-identification,open:rg-cover-holonomy-recovery}, and the regular quotient of the passive
coordinate action is \Cref{open:gen-gauge-quotient}.

\section{Base curves and contexts}
\label{sec:geo-base-curves}

Away from \(c_*\) the contextual base \(\mathcal C\) of \Cref{def:geo-context-base} is a smooth
manifold whose curves order contexts. A connection on a bundle over \(\mathcal C\) defines
transport in that bundle along base curves; unless a connection on \(T\mathcal C\) is separately
chosen, it does not transport tangent vectors of the base. A finite design
(\Cref{def:prob-finite-design}) is a finite subset of \(\mathcal C\), not a continuum limit or a
discretization of it. A base curve does not by itself represent the evolution of an agent. Change of
a belief or model at one context is a curve inside one associated fiber, simultaneous change over all
contexts is a curve in a declared space of sections, and a trajectory whose projection moves in
\(\mathcal C\) needs the horizontal and vertical splitting of a chosen connection; the three curve
types are defined in \Cref{def:hist-curve-types}.

\section{Relative frames and the two-atlas apparatus}
\label{sec:geo-relative-frames}

The main manuscript uses one frame per agent, a gauge choice made because no energy, kernel, or
evaluator at \(c_*\) reads a relative frame, and it records at one fiber how two atlases reduce to one
(\Cref{prop:geo-common-frame-gauge}). Belief and model frames may instead be chosen as two local
sections of the same principal bundle over the supports \(\mathcal C_i\). In this supplement a passive
re-expression of either frame is written with the convention of the main manuscript,
\(u'=u\cdot g^{-1}\), separately in each channel.

\propositionheading{Relative principal frame over a support}{prop:geo-relative-frame}
For each agent \(i\) let \(u_i^b,u_i^m:\mathcal C_i\to\mathscr P_G\) be smooth local sections of the
principal bundle. (i) The relative frame \(h_i\) of \eqref{eq:geo-relative-frame}, taken at every
point of \(\mathcal C_i\), is a smooth map \(h_i:\mathcal C_i\to G\), and its existence uses only
\(\mathscr P_G\), not the representations \(\rho_b,\rho_m\). (ii) Under independent smooth passive
re-expressions
\begin{equation}
u_i^{b\prime}=u_i^b\cdot(g_i^b)^{-1},
\qquad
u_i^{m\prime}=u_i^m\cdot(g_i^m)^{-1},
\qquad
g_i^b,g_i^m:\mathcal C_i\longrightarrow G,
\label{eq:geo-two-atlas-reframing}
\end{equation}
associated coordinates change by \(\beta_x'=\widehat\rho_x(g_i^x)\beta_x\) and \(h_i\) obeys
\eqref{eq:geo-relative-frame-law} pointwise. (iii) At \(c\in\mathcal C_i\cap\mathcal C_j\) the
comparisons \(T_{ij}^x\in G\) defined by \(u_j^x(c)=u_i^x(c)\cdot T_{ij}^x\) transform by
\(T_{ij}^{x\prime}=g_i^xT_{ij}^x(g_j^x)^{-1}\) and satisfy \eqref{eq:geo-frame-comparison-relation}.
(iv) The represented pair \((\rho_b(h_i),\rho_m(h_i))\) determines \(h_i\) only modulo
\(\ker\rho_b\cap\ker\rho_m\). A basis of an associated vector fiber that is not induced by a
section of \(\mathscr P_G\) need not be related to an induced basis by any element of \(G\): for
\(G=\mathrm{SO}(2)\) acting on \(\R^2\) by its defining representation, the basis obtained from an
induced one by \(\operatorname{diag}(2,1)\) is not.
\status{ESTABLISHED}

\noindent\emph{Proof.}
(i) Pointwise existence and uniqueness is \Cref{prop:geo-common-frame-gauge}; in a local
trivialization in which \(u_i^x(c)=(c,U^x(c))\) with smooth \(U^x\), \(h_i=(U^b)^{-1}U^m\) is smooth.
(ii) The quotient convention \eqref{eq:geo-quotient-convention} gives
\([u,\beta]=[u'\cdot g,\beta]=[u',\widehat\rho_x(g)\beta]\), and
\(u_i^{m\prime}=u_i^bh_i(g_i^m)^{-1}=u_i^{b\prime}\cdot g_i^bh_i(g_i^m)^{-1}\) with uniqueness in
the torsor gives the law. (iii) From
\(u_j^{x\prime}=u_i^xT_{ij}^x(g_j^x)^{-1}=u_i^{x\prime}\cdot g_i^xT_{ij}^x(g_j^x)^{-1}\) follows the
transformation law, and the two expressions \(u_j^m=u_j^bh_j=u_i^bT_{ij}^bh_j\) and
\(u_j^m=u_i^mT_{ij}^m=u_i^bh_iT_{ij}^m\) give the relation. (iv)
\(\rho_x(h)=\rho_x(h')\) exactly when \(h^{-1}h'\in\ker\rho_x\). An induced basis is carried to
another induced basis by \(\rho_b(k)\in\mathrm{SO}(2)\), and \(\operatorname{diag}(2,1)\notin
\mathrm{SO}(2)\). \hfill$\square$

The relative frame compares two principal sections even when \(\rho_b\) and \(\rho_m\) are
inequivalent or act on fibers of different dimensions. It produces two represented operators,
\(\widehat\rho_b(h_i)\) on \(\mathcal B_b\) and \(\widehat\rho_m(h_i)\) on \(\mathcal B_m\), and no
operator from \(\mathcal B_b\) to \(\mathcal B_m\); a comparison of channel points still needs a
cross map of \Cref{def:geo-cross-morphisms}. At every level \(\ell\) of a scale diagram the same
construction applies: on a common trivializing neighborhood \(\mathcal U_{\ell,i}\) of the level
base, two sections determine
\begin{equation}
u_{\ell,i}^m=u_{\ell,i}^b\cdot h_{\ell,i},
\qquad
h_{\ell,i}:\mathcal U_{\ell,i}\longrightarrow G_\ell,
\label{eq:rg-relative-frame}
\end{equation}
with the law \eqref{eq:geo-relative-frame-law} under independent rechoices, by
\Cref{prop:geo-relative-frame} at level \(\ell\). The common-frame
convention of the main manuscript is \(h_i=e_G\) with \(g_i^b=g_i^m\).

The active gauge group is not enlarged by the second atlas. At one context the argument of
\Cref{prop:geo-active-image-point} applies with an \(h_i\)-twist in the model channel.

\propositionheading{Active gauge image under independent atlases}{prop:geo-two-atlas-active-image}
Fix \(c\in\mathcal C\) with \(V_c=\{i\in V:c\in\mathcal C_i\}\) nonempty, and evaluate \(h_i\) and
\(T_{ij}^x\) at \(c\). For a right-\(G\)-equivariant bijection \(F_c\) of the fiber
\((\mathscr P_G)_c\) define \(\varphi_i^x\in G\) by
\(F_c(u_i^x(c))=u_i^x(c)\cdot\varphi_i^x\). Then
\begin{equation}
\varphi_j^x=(T_{ij}^x)^{-1}\varphi_i^xT_{ij}^x,
\qquad
\varphi_i^m=h_i^{-1}\varphi_i^bh_i,
\qquad i,j\in V_c,
\label{eq:geo-two-atlas-active-relations}
\end{equation}
the second being \eqref{eq:geo-diagonal-gauge-functions}, and
\(F_c\mapsto(\varphi_i^b,\varphi_i^m)_{i\in V_c}\) is a group isomorphism from
\(\operatorname{Aut}_G((\mathscr P_G)_c)\) onto the subgroup
\(\mathcal A_c\subseteq\prod_{i\in V_c}(G\times G)\) cut out by
\eqref{eq:geo-two-atlas-active-relations}. For each \(i_0\in V_c\) the projection onto
\(\varphi_{i_0}^b\) is an isomorphism \(\mathcal A_c\cong G\). For nontrivial \(G\),
\(\mathcal A_c\) is a proper subgroup of the passive product \(\prod_{i\in V_c}(G\times G)\), and,
as soon as \(|V_c|\geq2\), a proper subgroup of the product
\(\prod_{i\in V_c}\{(k,h_i^{-1}kh_i):k\in G\}\) of pointwise twisted diagonals. The local functions
at \(c\) of every gauge automorphism of \(\mathscr P_G\) form one element of \(\mathcal A_c\).
\status{ESTABLISHED}

\noindent\emph{Proof.}
Equivariance gives
\(F_c(u_j^x)=F_c(u_i^xT_{ij}^x)=u_i^x\varphi_i^xT_{ij}^x=u_j^x(T_{ij}^x)^{-1}\varphi_i^xT_{ij}^x\),
and uniqueness in the torsor gives the first relation; the second follows from
\(F_c(u_i^m)=F_c(u_i^bh_i)=u_i^b\varphi_i^bh_i=u_i^mh_i^{-1}\varphi_i^bh_i\). Since
\(F_c(F_c'(u))=u\cdot\varphi\varphi'\), the map is a homomorphism, and it is injective because
\(\varphi_{i_0}^b\) determines \(F_c\) on the whole torsor. Given an element of \(\mathcal A_c\), the
map \(u_{i_0}^b(c)\cdot g\mapsto u_{i_0}^b(c)\cdot\varphi_{i_0}^bg\) is an equivariant bijection
whose local functions are the given ones, and every \(\varphi_{i_0}^b\in G\) occurs, which proves
\(\mathcal A_c\cong G\). For \(k\neq e_G\), the element with \(\varphi_{i}^b=k\), \(\varphi_i^m=h_i^{-1}
kh_i\) and \(\varphi_j^b=e_G=\varphi_j^m\) for \(j\neq i\) lies in the product of twisted diagonals
but violates the first relation of \eqref{eq:geo-two-atlas-active-relations}; the passive product is
larger still. A gauge automorphism restricts to an equivariant bijection of each fiber.
\hfill$\square$

Whether a given element of \(\mathcal A_c\) is the value of a global gauge automorphism is a
separate question about the base.

\propositionheading{Realization of fiber automorphisms by gauge automorphisms}{prop:geo-fiber-automorphism-realization}
Let \(c_1,\ldots,c_M\in\mathcal C\) be distinct and let
\(F_d\in\operatorname{Aut}_G((\mathscr P_G)_{c_d})\) for \(d=1,\ldots,M\). If \(G\) is connected,
there is a smooth gauge automorphism \(F\) of \(\mathscr P_G\) over \(\operatorname{id}_{\mathcal C}\)
that restricts to \(F_d\) at every \(c_d\). For disconnected \(G\) this can fail: gluing
\([0,1]\times S_3\) by \((1,g)\sim(0,\mathring g g)\) for a \(3\)-cycle \(\mathring g\) gives a principal
\(S_3\)-bundle over the circle, and the fiber automorphism \((0,g)\mapsto(0,\lambda g)\) with
\(\lambda\) a transposition is the restriction of no gauge automorphism.
\status{ESTABLISHED}

\noindent\emph{Proof.}
Choose pairwise disjoint neighborhoods of the \(c_d\) carrying local sections \(\sigma_d\), define
\(\kappa_d\in G\) by \(F_d(\sigma_d(c_d))=\sigma_d(c_d)\cdot\kappa_d\), and write
\(\kappa_d=\exp(\xi_1)\cdots\exp(\xi_r)\) with \(\xi_1,\ldots,\xi_r\in\mathfrak g\), which is
possible because a connected Lie group is generated by the image of its exponential map. Put
\(\gamma_d(t)=\exp(t\xi_1)\cdots\exp(t\xi_r)\). With smooth bump functions \(\chi_d\) of compact
support in those neighborhoods and \(\chi_d(c_d)=1\), the map
\(F(\sigma_d(y)\cdot g)=\sigma_d(y)\cdot\gamma_d(\chi_d(y))g\) there and \(F=\operatorname{id}\)
elsewhere is a smooth equivariant automorphism restricting to \(F_d\) at \(c_d\). For the \(S_3\)
bundle, a gauge automorphism acts on the strip by \((t,g)\mapsto(t,\kappa g)\) with \(\kappa\)
constant because \(S_3\) is discrete, and compatibility with the gluing forces
\(\kappa\mathring g=\mathring g\kappa\). The centralizer of a \(3\)-cycle in \(S_3\) is the cyclic group it
generates, which contains no transposition. \hfill$\square$

Graph links between fibers over different contexts are compared through link maps. Anchor every
vertex at a context \(c_i\in\mathcal C_i\), evaluate the frames and \(h_i\) there, and for an edge
copy \(e:j\to i\) of the interaction multigraph of \Cref{def:geo-graph-links} define
\(\mathsf t_e^x:(\mathscr P_G)_{c_j}\to(\mathscr P_G)_{c_i}\) by
\(\mathsf t_e^x(u_j^x\cdot g)=u_i^x\cdot\Theta_e^xg\).

\propositionheading{Identified links under independent atlases}{prop:geo-identified-links}
Each \(\mathsf t_e^x\) is a right-\(G\)-equivariant bijection that does not depend on the vertex
frames of its channel when the links transform by \(\Theta_e^{x\prime}=g_i^x\Theta_e^x(g_j^x)^{-1}\).
The two channels carry one link on \(e\), meaning \(\mathsf t_e^b=\mathsf t_e^m\), exactly when
\begin{equation}
\Theta_e^m=h_i^{-1}\Theta_e^bh_j,
\qquad e:j\to i,
\label{eq:geo-common-link-intrinsic}
\end{equation}
and this relation is invariant under independent re-expressions of both atlases. The literal
equality \(\Theta_e^b=\Theta_e^m\) is not invariant under them; under the common-frame
convention \(h_i=h_j=e_G\), it expresses \(\mathsf t_e^b=\mathsf t_e^m\), which is the form
\eqref{eq:geo-common-link} of the main manuscript. More generally, equal link coordinates
\(\Theta_e^b=\Theta_e^m=\Theta_e\) represent equal link maps exactly when
\(h_i^{-1}\Theta_e h_j=\Theta_e\). \status{ESTABLISHED}

\noindent\emph{Proof.}
Every point of \((\mathscr P_G)_{c_j}\) is \(u_j^x\cdot g\) for a unique \(g\), so
\(\mathsf t_e^x\) is well defined, equivariant, and invertible. After a rechoice,
\(\mathsf t_e^{x\prime}(u_j^{x\prime}\cdot g)=u_i^x(g_i^x)^{-1}g_i^x\Theta_e^x(g_j^x)^{-1}g
=u_i^x\Theta_e^x(g_j^x)^{-1}g=\mathsf t_e^x(u_j^x(g_j^x)^{-1}g)=\mathsf t_e^x(u_j^{x\prime}\cdot g)\).
Next, \(\mathsf t_e^b(u_j^m\cdot g)=\mathsf t_e^b(u_j^b\cdot h_jg)=u_i^b\Theta_e^bh_jg
=u_i^m\cdot h_i^{-1}\Theta_e^bh_jg\), which equals \(\mathsf t_e^m(u_j^m\cdot g)=u_i^m\Theta_e^mg\)
for every \(g\) exactly when \eqref{eq:geo-common-link-intrinsic} holds; invariance follows
because both link maps are frame independent. A rechoice of the model atlas alone with
\(g_j^m=e_G\) and \(g_i^m=k\neq e_G\) replaces \(\Theta_e^m\) by \(k\Theta_e^m\) and leaves
\(\Theta_e^b\) fixed, so it breaks the literal equality. \hfill$\square$

Cross maps are the remaining place where two atlases matter, because their local representatives
must be converted to one atlas before they can be read as sections of a map bundle.

\propositionheading{Cross maps under independent atlases}{prop:geo-map-bundle-conversion}
Let \(\Phi,\widetilde\Phi\) be cross morphisms (\Cref{def:geo-cross-morphisms}) with local
representatives \(\phi_i=\Phi_i^{bm}\) and \(\widetilde\phi_i\) in the frames \(u_i^b,u_i^m\). (i) Under
\eqref{eq:geo-two-atlas-reframing},
\(\phi_i'=\widehat\rho_m(g_i^m)\circ\phi_i\circ\widehat\rho_b(g_i^b)^{-1}\) and
\(\widetilde\phi_i'=\widehat\rho_b(g_i^b)\circ\widetilde\phi_i\circ\widehat\rho_m(g_i^m)^{-1}\).
(ii) On overlaps
\(\phi_i\circ\widehat\rho_b(T_{ij}^b)=\widehat\rho_m(T_{ij}^m)\circ\phi_j\) and
\(\widetilde\phi_i\circ\widehat\rho_m(T_{ij}^m)=\widehat\rho_b(T_{ij}^b)\circ\widetilde\phi_j\).
(iii) In the belief atlas the converted representatives
\begin{equation}
\check\phi_i:=\widehat\rho_m(h_i)\circ\phi_i,
\qquad
\check{\widetilde\phi}_i:=\widetilde\phi_i\circ\widehat\rho_m(h_i^{-1}),
\label{eq:geo-two-atlas-map-conversion}
\end{equation}
the first being \(\Phi_i^{\mathrm{conv}}\) of \eqref{eq:geo-map-bundle-frame-conversion}, satisfy
\(\check\phi_i\circ\widehat\rho_b(T_{ij}^b)=\widehat\rho_m(T_{ij}^b)\circ\check\phi_j\) and
\(\check{\widetilde\phi}_i\circ\widehat\rho_m(T_{ij}^b)=\widehat\rho_b(T_{ij}^b)\circ
\check{\widetilde\phi}_j\), so they are local representatives of sections of the associated map
bundles with actions \(\phi\mapsto\widehat\rho_m(g)\circ\phi\circ\widehat\rho_b(g)^{-1}\) and its
interchange. (iv) If \(\Phi\) is induced by an intertwiner \(\phi\) as in
\eqref{eq:geo-induced-cross-morphisms}, then \(\phi_i=\widehat\rho_m(h_i^{-1})\circ\phi\) and
\(\check\phi_i=\phi\). \status{ESTABLISHED}

\noindent\emph{Proof.}
A point \([u_i^b,\beta]\) is sent to \([u_i^m,\phi_i\beta]\). In new coordinates its input is
\(\widehat\rho_b(g_i^b)\beta\) and its output \(\widehat\rho_m(g_i^m)\phi_i\beta\), which gives (i).
Over \(c\in\mathcal C_i\cap\mathcal C_j\), \([u_j^b,\beta]=[u_i^b,\widehat\rho_b(T_{ij}^b)\beta]\) has
image \([u_j^m,\phi_j\beta]=[u_i^m,\widehat\rho_m(T_{ij}^m)\phi_j\beta]\) computed in \(j\) and
\([u_i^m,\phi_i\widehat\rho_b(T_{ij}^b)\beta]\) computed in \(i\), which gives (ii); the second map
is the same computation. For (iii),
\(\check\phi_i\widehat\rho_b(T_{ij}^b)=\widehat\rho_m(h_iT_{ij}^m)\phi_j
=\widehat\rho_m(T_{ij}^bh_j)\phi_j=\widehat\rho_m(T_{ij}^b)\check\phi_j\) by
\eqref{eq:geo-frame-comparison-relation}, and
\(\check{\widetilde\phi}_i\widehat\rho_m(T_{ij}^b)=\widetilde\phi_i\widehat\rho_m(T_{ij}^mh_j^{-1})
=\widehat\rho_b(T_{ij}^b)\check{\widetilde\phi}_j\). For (iv),
\([u_i^b,\beta]\mapsto[u_i^b,\phi\beta]=[u_i^mh_i^{-1},\phi\beta]=[u_i^m,\widehat\rho_m(h_i^{-1})\phi\beta]\).
\hfill$\square$

The compatible family \((\phi_i)\) is therefore a map-bundle section only after the conversion
\eqref{eq:geo-two-atlas-map-conversion}, and \(h_i\) does not replace the cross-fiber map.

\section{Connections, parallel transport, and covariant defects}
\label{sec:geo-three-transports}

\definitionheading{Two connections on the common principal bundle}{def:geo-connections}
Under \Cref{hyp:geo-smooth-tier}, choose principal connections
\begin{equation}
\omega_b,\omega_m\in\Omega^1(\mathscr P_G,\mathfrak g).
\label{eq:geo-principal-connections}
\end{equation}
They induce horizontal distributions \(H^b,H^m\) on \(\mathcal E_b,\mathcal E_m\), with horizontal
lift \(H_e^xX\) of \(X\in T_c\mathcal C\) at \(e\in(\mathcal E_x)_c\), and, for every
piecewise-smooth \(\gamma:[0,1]\to\mathcal C\), parallel transports
\begin{equation}
\Omega_\gamma:
(\mathcal E_b)_{\gamma(0)}\longrightarrow(\mathcal E_b)_{\gamma(1)},
\qquad
\widetilde\Omega_\gamma:
(\mathcal E_m)_{\gamma(0)}\longrightarrow(\mathcal E_m)_{\gamma(1)},
\label{eq:geo-base-parallel-transports}
\end{equation}
obtained by carrying the associated-law coordinate along the \(\omega_x\)-horizontal lift of
\(\gamma\) through a point of \((\mathscr P_G)_{\gamma(0)}\).

This is transport over the base curve, not transport of vectors in \(T\mathcal C\); a
tangent-vector transport needs a separately chosen affine connection on \(T\mathcal C\), unless
\(\mathscr P_G\) and its representation have been identified with the frame bundle of
\(\mathcal C\). The two connections need not agree. Their difference \(\omega_m-\omega_b\) is
horizontal and \(\operatorname{Ad}\)-equivariant, hence descends to an
\(\operatorname{Ad}(\mathscr P_G)\)-valued one-form on \(\mathcal C\). In an associated vector-bundle
realization the horizontal distributions define covariant derivatives \(\nabla^b,\nabla^m\); no
linear covariant derivative is assumed on a law fiber. The equal-connection specialization declares
\(\omega_b=\omega_m\); the principal holonomy is then common, while its actions on belief and model
laws remain \(\widehat\rho_b\) and \(\widehat\rho_m\). It constrains the connections only, so it is
separate from the common frame and from identified links, which constrain frames and link data.

\propositionheading{Frame laws of connections and transports}{prop:geo-connection-frame-laws}
Let \(A_i^x=(u_i^x)^*\omega_x\). In matrix notation for the group, the re-expressions
\eqref{eq:geo-two-atlas-reframing} give
\begin{align}
A_i^{b\prime}
&=\operatorname{Ad}_{g_i^b}A_i^b-dg_i^b(g_i^b)^{-1},
\label{eq:geo-local-connection-b}\\
A_i^{m\prime}
&=\operatorname{Ad}_{g_i^m}A_i^m-dg_i^m(g_i^m)^{-1}.
\label{eq:geo-local-connection-m}
\end{align}
If the endpoint coordinates of \(\gamma\) are re-expressed by \(g_0,g_1\) in the belief frames and by
\(\widetilde g_0,\widetilde g_1\) in the model frames, the represented transports obey
\begin{align}
\Omega_\gamma'
&=\widehat\rho_b(g_1)\circ\Omega_\gamma\circ\widehat\rho_b(g_0)^{-1},
\label{eq:geo-pt-gauge-b}\\
\widetilde\Omega_\gamma'
&=\widehat\rho_m(\widetilde g_1)\circ\widetilde\Omega_\gamma\circ\widehat\rho_m(\widetilde g_0)^{-1}.
\label{eq:geo-pt-gauge-m}
\end{align}
\status{ESTABLISHED}

\noindent\emph{Proof.}
In a trivialization with \(u(x)=(x,U(x))\), a principal connection has
\(u^*\omega=U^{-1}A_0U+U^{-1}dU\) for the local form \(A_0\) of the identity section. For
\(U'=Ug^{-1}\) this gives
\(gU^{-1}A_0Ug^{-1}+gU^{-1}dUg^{-1}+g\mathrm d(g^{-1})=\operatorname{Ad}_gA-(\mathrm dg)g^{-1}\), because
\(\mathrm d(g^{-1})=-g^{-1}(\mathrm dg)g^{-1}\). The transport is a coordinate-free map between fibers, and the
coordinates at each endpoint change by \(\beta'=\widehat\rho_x(g)\beta\). These are the standard laws
for the quotient convention \eqref{eq:geo-quotient-convention}
\citep{Kobayashi1963,Nakahara2003}. \hfill$\square$

The endpoint laws also show why a same-point frame comparison and transport along a nonconstant
curve are different objects: the first involves one fiber, the second two.

\definitionheading{Covariant defects}{def:geo-covariant-defects}
Let \(e\in\mathcal E_b\) and \(e'\in\mathcal E_m\) lie over \(c\), and let \(X\in T_c\mathcal C\). For
smooth associated fibers and smooth cross morphisms define the vertical-valued tensors
\begin{align}
(\mathcal D\Phi)_e(X)
&=T_e\Phi\bigl(H_e^bX\bigr)-H_{\Phi(e)}^mX,
\label{eq:geo-nonlinear-defect-phi}\\
(\mathcal D\widetilde\Phi)_{e'}(X)
&=T_{e'}\widetilde\Phi\bigl(H_{e'}^mX\bigr)-H_{\widetilde\Phi(e')}^bX.
\label{eq:geo-nonlinear-defect-tildephi}
\end{align}
The cross map is parallel when its defect vanishes identically.

\propositionheading{Covariant defects and parallelness}{prop:geo-covariant-defect-parallelness}
(i) Both defects take vertical values. (ii) \(\mathcal D\Phi=0\) and \(\mathcal D\widetilde\Phi=0\)
hold exactly when every piecewise-smooth base curve satisfies
\begin{align}
\widetilde\Omega_\gamma\circ\Phi_{\gamma(0)}
&=\Phi_{\gamma(1)}\circ\Omega_\gamma,
\label{eq:geo-phi-parallelness}\\
\Omega_\gamma\circ\widetilde\Phi_{\gamma(0)}
&=\widetilde\Phi_{\gamma(1)}\circ\widetilde\Omega_\gamma,
\label{eq:geo-tildephi-parallelness}
\end{align}
respectively. (iii) In a fiberwise-linear associated-vector-bundle specialization with
representation spaces \(\mathbb V_b,\mathbb V_m\), the defects reduce to
\begin{equation}
D\Phi=\nabla^m\circ\Phi-\Phi\circ\nabla^b,
\qquad
D\widetilde\Phi=\nabla^b\circ\widetilde\Phi-\widetilde\Phi\circ\nabla^m,
\label{eq:geo-linear-covariant-defects}
\end{equation}
the old-to-new coordinate maps induced by \eqref{eq:geo-two-atlas-reframing} are
\begin{equation}
R_i^b=\rho_b(g_i^b):\mathbb V_b\longrightarrow\mathbb V_b,
\qquad
R_i^m=\rho_m(g_i^m):\mathbb V_m\longrightarrow\mathbb V_m,
\label{eq:geo-linear-coordinate-maps}
\end{equation}
and the defect matrices obey
\begin{equation}
(D\Phi)_i'=R_i^m(D\Phi)_i(R_i^b)^{-1},
\qquad
(D\widetilde\Phi)_i'=R_i^b(D\widetilde\Phi)_i(R_i^m)^{-1}.
\label{eq:geo-defect-gauge-laws}
\end{equation}
\status{ESTABLISHED}

\noindent\emph{Proof.}
(i) The two terms of each defect project to the same \(X\). (ii) Vanishing defect says that the
image of every horizontal lift is horizontal; integration along \(\gamma\) gives the two
commutative squares, and differentiation at the initial point gives the converse. (iii) In the
linear specialization \(T\Phi\) of a horizontal lift minus the horizontal lift of the image is the
covariant derivative of the Hom-bundle section; the coordinate maps are \(R=\rho_x(g)\) by
\(\beta'=\widehat\rho_x(g)\beta\); and with \(\nabla^{x\prime}=R^x\nabla^x(R^x)^{-1}\) and
\(\Phi'=R^m\Phi(R^b)^{-1}\) one gets
\(\nabla^{m\prime}\Phi'-\Phi'\nabla^{b\prime}=R^m(\nabla^m\Phi-\Phi\nabla^b)(R^b)^{-1}\).
\hfill$\square$

The coordinate maps \(R_i^x\) are invertible re-expressions within one channel, not cross-channel
morphisms or coarse restriction operators.

\section{The \v{C}ech picture and bundle topology}
\label{sec:geo-cech}

Let \(\mathcal C_0=\bigcup_i\mathcal C_i\), and let \(\mathcal N\) be the nerve of this open cover,
\begin{equation}
\{i_0,\ldots,i_r\}\in\mathcal N
\quad\Longleftrightarrow\quad
\mathcal C_{i_0\cdots i_r}=\bigcap_{k=0}^r\mathcal C_{i_k}\neq\varnothing.
\label{eq:geo-nerve}
\end{equation}
Write \(\underline G\) for the sheaf of smooth \(G\)-valued maps.

\propositionheading{Transition cocycles and the principal class}{prop:geo-cech-class}
(i) For either atlas \(x\in\{b,m\}\), the comparisons \(T_{ij}^x:\mathcal C_{ij}\to G\) obey
\begin{equation}
T_{ii}^x=e_G,
\qquad
T_{ji}^x=(T_{ij}^x)^{-1},
\qquad
T_{ij}^xT_{jl}^x=T_{il}^x.
\label{eq:geo-cech-cocycle}
\end{equation}
(ii) The two cocycles are cohomologous and define one nonabelian \v{C}ech class,
\begin{equation}
[T^b]=[T^m]=[\mathscr P_G|_{\mathcal C_0}]
\in\check H^1(\{\mathcal C_i\};\underline G).
\label{eq:geo-cech-class}
\end{equation}
(iii) The class is distinguished exactly when \(\mathscr P_G|_{\mathcal C_0}\) is trivial,
equivalently when, for either atlas, there are smooth \(U_i^x:\mathcal C_i\to G\) with
\begin{equation}
T_{ij}^x=U_i^x(U_j^x)^{-1}.
\label{eq:geo-coboundary-form}
\end{equation}
Over a single context every principal bundle is trivial and every class is distinguished.
(iv) The single principal class does not make the associated bundles isomorphic: for the
nontrivial principal \(\{\pm1\}\)-bundle over the circle, the trivial and the sign representations
on \(\R\) give the trivial line bundle and the M\"obius line bundle, which are not isomorphic.
(v) A global section of an associated bundle need not trivialize \(\mathscr P_G\): if a fiber
contains a \(G\)-fixed point \(\beta_0\), then \(c\mapsto[p,\beta_0]\) is a global section for every
\(\mathscr P_G\). \status{ESTABLISHED}

\noindent\emph{Proof.}
(i) On a triple overlap, composing the unique change from \(l\) to \(j\) with the unique change from
\(j\) to \(i\) gives the unique change from \(l\) to \(i\). (ii) Relation
\eqref{eq:geo-frame-comparison-relation} is a coboundary equivalence of the two cocycles. (iii) If
\eqref{eq:geo-coboundary-form} holds, the rechoice with \(g_i^x=(U_i^x)^{-1}\) in
\eqref{eq:geo-two-atlas-reframing} makes every transition the identity, so the sections glue to a
global section, which trivializes a principal bundle; conversely a global trivialization supplies the
\(U_i^x\). This is the global-section characterization of trivial principal bundles
\citep{Steenrod1951,Husemoller1994,Kobayashi1963}. A single context is a point, over which every
bundle is trivial. (iv) A nowhere-vanishing continuous section of the M\"obius bundle would be a
continuous \(\sigma:[0,1]\to\R\setminus\{0\}\) with \(\sigma(1)=-\sigma(0)\), which the
intermediate value theorem excludes, while the trivial line bundle has the section \(1\). (v)
\([p\cdot g,\beta_0]=[p,g\cdot\beta_0]=[p,\beta_0]\), so the value does not depend on the
representative. \hfill$\square$

\hypothesisheading{Common principal trivialization}{hyp:geo-common-trivializations}
Whenever global principal-frame coordinates are used, \(\mathscr P_G|_{\mathcal C_0}\) is trivial
and one reference section \(\sigma_0\) is chosen, so that
\begin{equation}
u_i^b=\sigma_0\cdot(U_i^b)^{-1},
\qquad
u_i^m=\sigma_0\cdot(U_i^m)^{-1},
\qquad
h_i=U_i^b(U_i^m)^{-1}.
\label{eq:geo-frame-field}
\end{equation}
The hypothesis excludes nontrivial restrictions \(\mathscr P_G|_{\mathcal C_0}\). It is one
topological hypothesis, not two channelwise ones. At \(c_*\) the main manuscript needs no such
hypothesis, because a reference point of the single fiber supplies the potentials
(\Cref{def:geo-principal-systems}).

\subsection{Independent channel frames in the Gaussian realization}
\label{sec:gen-gaussian-covariance}

Supplement G realizes the directed law with Gaussian kernels in the common frame. With independent
atlases, invoke \Cref{hyp:geo-common-trivializations} on a declared region \(\mathcal C_0\)
containing the contexts under study, restrict each frame to \(\mathcal C_i\cap\mathcal C_0\), and
write the independent frame coordinates as
\begin{equation}
U_i^b,U_i^m:\mathcal C_i\cap\mathcal C_0\longrightarrow G,
\qquad
u_i^x=\sigma_0\cdot(U_i^x)^{-1},
\label{eq:gen-frames}
\end{equation}
so that \(\sigma_0\) is common only on the trivialized region and the relative frame is
\begin{equation}
h_i=U_i^b(U_i^m)^{-1},
\qquad
u_i^m=u_i^bh_i.
\label{eq:gen-relative-principal-frame}
\end{equation}
Its represented actions \(\rho_b(h_i)\) and \(\rho_m(h_i)\) live in different endomorphism spaces
and define neither the bridge matrix of the Gaussian model nor either cross morphism. An independent
pair \((g_i^b,g_i^m)\in G\times G\) has represented matrices \(R_i^b=\rho_b(g_i^b)\) and
\(R_i^m=\rho_m(g_i^m)\); the pair records two trivializations of bundles associated to one
\(\mathscr P_G\) and does not replace the structure group by \(G\times G\).

\remarkheading{Gaussian kernel covariance under independent frames}{rem:gen-independent-frame-covariance}
The transformation laws of the Gaussian directed parameters in Supplement G,
\cref{eq:gen-gauge-selfmaps,eq:gen-gauge-bridge-offsets,eq:gen-gauge-covariances,eq:gen-gauge-precision,eq:gen-gauge-root-model,eq:gen-gauge-root-state,eq:gen-gauge-observation},
hold verbatim with \(R_i^b=\rho_b(g_i^b)\) and \(R_i^m=\rho_m(g_i^m)\) for independent
\(g_i^b,g_i^m\) and links transforming by \(\Theta_e^{x\prime}=g_i^x\Theta_e^x(g_j^x)^{-1}\), and
the kernel covariance of \Cref{prop:gen-lg-kernel-covariance} holds under every independent
channel-frame re-expression. \status{ESTABLISHED}

\noindent\emph{Proof.}
Each law and each step of the proof of \Cref{prop:gen-lg-kernel-covariance} uses the belief matrix
only on belief coordinates and the model matrix only on model coordinates: with
\(F'=R_i^bF(R_i^b)^{-1}\), \(\rho_b(\Theta_e^{b\prime})=R_i^b\rho_b(\Theta_e^b)(R_j^b)^{-1}\),
\(B'=R_i^bB(R_i^m)^{-1}\), and \(\eta'=R_i^b\eta\), the transformed belief mean at transformed inputs
is \(R_i^b\) times the original mean, and the covariance congruence makes each transformed kernel the
receiving-frame pushforward. No step uses a relation between \(R_i^b\) and \(R_i^m\).
\hfill$\square$

\section{Graph links and base transport}
\label{sec:geo-graph-base-transport}

The links of \Cref{def:geo-graph-links} are declared independently of \(\mathscr P_G\) and of any
connection. Relating them to base transport is a separate declaration.

\hypothesisheading{Graph links realized by base transport}{hyp:geo-graph-base-transport}
Contexts \(c_i\in\mathcal C_i\) are chosen, and for every edge copy \(e:j\to i\) a piecewise-smooth
curve \(\gamma_e\) from \(c_j\) to \(c_i\) is assigned, with \(\gamma_{\bar e}\) the reverse of
\(\gamma_e\). In the chosen endpoint frames,
\begin{equation}
\widehat\rho_b(\Theta_e^b)=\Omega_{\gamma_e},
\qquad
\widehat\rho_m(\Theta_e^m)=\widetilde\Omega_{\gamma_e}.
\label{eq:geo-link-pt-hypothesis}
\end{equation}
The hypothesis excludes links that are not represented transports of the chosen connections along
the assigned curves.

\propositionheading{Loop products under base-transport realization}{prop:geo-graph-base-loop-transport}
Under \Cref{hyp:geo-graph-base-transport}, for every closed edge-labeled walk
\(p=(e_0,\ldots,e_{r-1})\) with \(e_k:i_{k+1}\to i_k\) and \(i_r=i_0\),
\[
\widehat\rho_b\bigl(\mathrm{Hol}^b(p)\bigr)=\Omega_{\gamma_p},
\qquad
\widehat\rho_m\bigl(\mathrm{Hol}^m(p)\bigr)=\widetilde\Omega_{\gamma_p},
\]
where \(\gamma_p\) is the base loop at \(c_{i_0}\) that traverses
\(\gamma_{e_{r-1}},\gamma_{e_{r-2}},\ldots,\gamma_{e_0}\) in this order. \status{ESTABLISHED}

\noindent\emph{Proof.}
The representation is a homomorphism, so
\(\widehat\rho_x(\Theta_{e_0}^x\cdots\Theta_{e_{r-1}}^x)\) is the composite of the represented links
with \(\widehat\rho_x(\Theta_{e_{r-1}}^x)\) applied first; consecutive factors are expressed in the
same frame at their common endpoint, and parallel transport along a concatenation is the composite
of the transports. \hfill$\square$

Without the curve assignment and the equalities \eqref{eq:geo-link-pt-hypothesis}, graph holonomy
and base-connection holonomy are unrelated, and even an exact assignment concerns connection
holonomy, not topological nontriviality of \(\mathscr P_G\). The hypothesis is satisfiable by
declaration, by setting each link equal to a chosen base transport, and no result of this supplement
realizes it for independently specified links; the cover construction of
\Cref{ch:neighborhood-parents} declines that identification, and the operational record of a
declared base holonomy is \Cref{thm:rg-cover-operational-holonomy}.

Whether independently specified links \(\Theta_e^b,\Theta_e^m\) are realized by principal
connections and base curves satisfying \eqref{eq:geo-link-pt-hypothesis}, and whether a connection
can be recovered from the laws of a declared class of population records, are
\Cref{open:rg-cover-graph-base-identification,open:rg-cover-holonomy-recovery}; what holds for
declared realizations is \Cref{prop:geo-graph-base-loop-transport}.

The single-point theory cannot separate the following conditions, because over one context every
bundle is trivial and every connection vanishes. Over a base they separate completely.

\propositionheading{Four independent notions of flatness}{prop:geo-four-notions}
Consider (a) triviality of \(\mathscr P_G\) over a region, equivalently a distinguished class in
\Cref{prop:geo-cech-class}; (b) graph flatness of the links of one channel in the sense of
\Cref{prop:geo-trivializing-criterion}; (c) flatness of a principal connection, meaning vanishing
curvature; and (d) triviality of the holonomy of that connection around one specified base loop.
Then (a) implies neither (c) nor (d), (c) implies neither (a) nor (d), and (d) does not imply (c);
and (b) is implied by none of (a), (c), (d) and implies none of them. Explicitly, for the trivial
\(\mathrm U(1)\)-bundle with connection form \(\mathrm i\alpha\) for a real one-form \(\alpha\),
whose holonomy around a loop \(\gamma\) is \(e^{\pm\mathrm i\oint_\gamma\alpha}\) with the sign fixed
by the orientation convention: \(\alpha=x\mathrm dy\) on \(\R^2\) is not flat and has trivial holonomy
around the circle of radius \(\sqrt2\); \(\alpha=\lambda(x\mathrm dy-y\mathrm dx)/(x^2+y^2)\) on
\(\R^2\setminus\{0\}\) with \(\lambda\notin\mathbb Z\) is flat with nontrivial holonomy around the
unit circle; and the principal \(S_3\)-bundle of \Cref{prop:geo-fiber-automorphism-realization} is
nontrivial while its unique connection is flat, with holonomy \(\mathring g\) around the base circle.
\status{ESTABLISHED}

\noindent\emph{Proof.}
For \(\alpha=x\mathrm dy\), \(\mathrm d\alpha=\mathrm dx\wedge\mathrm dy\neq0\), while \(\oint x\mathrm dy\) over the positively oriented
circle of radius \(\sqrt2\) is its enclosed area \(2\pi\), so the holonomy is trivial; this proves
that (d) does not imply (c) and, since the bundle is trivial, that (a) does not imply (c). For
\(\alpha=\lambda(x\mathrm dy-y\mathrm dx)/(x^2+y^2)\), \(\mathrm d\alpha=0\) on the punctured plane and
\(\oint\alpha=2\pi\lambda\notin2\pi\mathbb Z\) around the unit circle, so the connection is flat on
a trivial bundle with nontrivial holonomy; hence (a) does not imply (d) and (c) does not imply (d).
For the \(S_3\)-bundle, a continuous section of the strip is a constant \((t,s)\) because \(S_3\) is
discrete, and the gluing requires \(\mathring g s=s\), which fails, so the bundle is nontrivial; the Lie
algebra of \(S_3\) is zero, so its only connection has horizontal space equal to the whole tangent
space and zero curvature, and its holonomy around the circle is \(\mathring g\neq e\). Hence (c) does not
imply (a). Finally the links are declared independently of the bundle and the connection: the
all-identity links are graph flat over every bundle and connection, and links on a triangle whose loop
product differs from \(e_G\) are not graph flat over the trivial bundle with the zero connection,
which satisfies (a), (c), and (d). \hfill$\square$

\section{Optional product-gauge extension}
\label{sec:geo-product-extension}

\definitionheading{Product-gauge extension}{def:geo-product-extension}
An application that requires independently nonisomorphic principal bundles or two independent gauge
symmetries replaces \(\mathscr P_G\) by
\begin{equation}
\pi_b:P_b\to\mathcal C,
\qquad
\pi_m:P_m\to\mathcal C,
\qquad
Q=P_b\times_{\mathcal C}P_m,
\label{eq:geo-product-principal-extension}
\end{equation}
where \(Q\) is a principal \(G_b\times G_m\)-bundle and each associated bundle is formed with its own
factor. In a category admitting the mapping object, a belief-to-model morphism is a section of
\begin{equation}
Q\times_{G_b\times G_m}\operatorname{Map}(\mathcal B_b,\mathcal B_m),
\qquad
(g_b,g_m)\cdot\phi
=\widehat\rho_m(g_m)\circ\phi\circ\widehat\rho_b(g_b)^{-1}.
\label{eq:geo-product-cross-map-bundle}
\end{equation}
At every level \(\ell\) of a scale diagram the same construction uses
\(P_{\ell,b},P_{\ell,m}\to\mathcal C_\ell\) and \(Q_\ell=P_{\ell,b}\times_{\mathcal C_\ell}P_{\ell,m}\).

\propositionheading{Comparison data for the product extension}{prop:geo-product-extension-data}
(i) If \(G_b=G_m=G\) and \(B:P_b\to P_m\) is a principal \(G\)-bundle isomorphism over
\(\operatorname{id}_{\mathcal C}\), local sections \(u_i^b\) of \(P_b\) and \(u_i^m\) of \(P_m\)
determine a unique relative frame \(h_i\) with \(u_i^m=B(u_i^b)\cdot h_i\). (ii) If
\(\kappa:G_b\to G_m\) is a homomorphism and \(B:P_b\to P_m\) a bundle map over
\(\operatorname{id}_{\mathcal C}\) with \(B(p\cdot g_b)=B(p)\cdot\kappa(g_b)\), then every fiber map
\(\phi:\mathcal B_b\to\mathcal B_m\) with
\(\phi\circ\widehat\rho_b(g_b)=\widehat\rho_m(\kappa(g_b))\circ\phi\) induces the cross morphism
\([p,\beta]\mapsto[B(p),\phi(\beta)]\). (iii) A constant fiber map \(\phi\) defines a section of
\eqref{eq:geo-product-cross-map-bundle} in any chosen trivialization of \(Q\). This construction
is independent of the choice of trivialization exactly when
\((g_b,g_m)\cdot\phi=\phi\) for every \((g_b,g_m)\in G_b\times G_m\). \status{ESTABLISHED}

\noindent\emph{Proof.}
(i) \(B(u_i^b)\) and \(u_i^m\) are local sections of \(P_m\), so the torsor argument of
\Cref{prop:geo-relative-frame} applies. (ii) The representatives \((p\cdot g_b,\beta)\) and
\((p,\widehat\rho_b(g_b)\beta)\) have images \([B(p)\kappa(g_b),\phi\beta]\) and
\([B(p),\phi\widehat\rho_b(g_b)\beta]=[B(p),\widehat\rho_m(\kappa(g_b))\phi\beta]\), which agree.
(iii) The class \([q,\phi]\) is independent of the chosen principal frame over its base point exactly when
\([q\cdot(g_b,g_m),\phi]=[q,(g_b,g_m)\cdot\phi]\) equals \([q,\phi]\). \hfill$\square$

Without such comparison data the product bundle supplies neither a relative frame valued in one
group nor either cross morphism, and the invariance required in (iii) under independently varying
\((g_b,g_m)\) is restrictive. The extension carries additional topology and an independent product
gauge freedom, so it is invoked only when those degrees of freedom are intended, not because belief
and model use different representations of one group.

\section{Finite designs, context-indexed sections, and sampling compatibility}
\label{sec:prob-finite-designs}

\definitionheading{Finite design and incidence}{def:prob-finite-design}
A \emph{design} is a finite set of distinct contexts
\begin{equation}
D=\{c_d\}_{d=1}^{M}\subseteq\mathcal C,
\label{eq:prob-design}
\end{equation}
indexed by \(d\). The agent--design incidence set is \(I_D=\{(i,d):c_d\in\mathcal C_i\}\), with
\(D_i=\{d:(i,d)\in I_D\}\) and \(V_d=\{i:(i,d)\in I_D\}\); design points with no participating agent
are omitted, so \(V_d\neq\varnothing\). At each \((i,d)\in I_D\) measurable state and model fibers
\((\mathsf K_{i,d},\mathscr K_{i,d})\) and \((\mathsf M_{i,d},\mathscr M_{i,d})\) are declared, the
latent space is
\(\mathsf Y_D=\prod_{(i,d)\in I_D}(\mathsf K_{i,d}\times\mathsf M_{i,d})\) with its product
sigma-algebra, and \(\mathsf O_D\) is the observation space of the records declared on the design.
When \(D\subseteq\bigcap_{i\in V}\mathcal C_i\) the incidence is rectangular,
\(I_D=V\times\{1,\ldots,M\}\). The construction of the main manuscript is the case \(D=\{c_*\}\).

A sparse design is handled by indexing over declared incidences, never by writing a product over all
\((i,d)\) and deleting undefined coordinates, which would silently change the sample space. Missing
or aggregated observations are likewise represented by a declared observation map or kernel, which
changes \(\mathsf O_D\) and not \(\mathsf Y_D\). The recognition marginals at an incidence,
\(q_{i,d}^{b;o,X}\) and \(q_{i,d}^{m;o,X}\), are the marginals of \(\mathbb Q_{V,o,X}\) on
\(\mathsf K_{i,d}\) and \(\mathsf M_{i,d}\), as in \Cref{def:prob-recognition-marginals}.

\hypothesisheading{Measurable context-indexed section response}{hyp:prob-context-section-response}
For every agent \(i\), and for the structural space \(\mathsf X\) of
\Cref{def:prob-structural-kernel-signatures}, measurable spaces \((\mathfrak S_i^b,\mathscr S_i^b)\)
and \((\mathfrak S_i^m,\mathscr S_i^m)\) are declared whose underlying sets satisfy
\(\mathfrak S_i^x\subseteq\Gamma(\mathcal C_i,\mathcal E_x|_{\mathcal C_i})\), such that each
evaluation map \(\operatorname{eval}_{i,c}^x:\mathfrak S_i^x\to(\mathcal E_x)_c\), \(s\mapsto s(c)\),
is measurable, together with a measurable response map
\begin{equation}
\operatorname{Resp}_i:
(\mathsf O_D\times\mathsf X,\mathscr O_D\otimes\mathscr X)
\longrightarrow
(\mathfrak S_i^b\times\mathfrak S_i^m,
 \mathscr S_i^b\otimes\mathscr S_i^m),
\qquad
(o,X)\longmapsto(q_i^{b;o,X},q_i^{m;o,X}).
\label{eq:prob-context-section-response}
\end{equation}
At every \((i,d)\in I_D\) the fixed fiber identifications send the evaluated values into
\(\mathcal P(\mathsf K_{i,d},\mathscr K_{i,d})\) and \(\mathcal P(\mathsf M_{i,d},\mathscr M_{i,d})\),
so design-point evaluations are measurable functions of \((o,X)\). For fixed \((o,X)\) these sections
occupy the section slots of \eqref{eq:geo-agent-sections}. The hypothesis excludes responses whose
section values are not measurable in the observation and the structural data.

Evaluation and section evaluation have different codomains. For a sample-field realization
\(y=((k_i,m_i))_{i\in V}\) on the admitted incidences, the evaluation
\(\bigl((k_i(c_d),m_i(c_d))_{i\in V_d}\bigr)_{d=1}^{M}\) is a point of \(\mathsf Y_D\), whereas
\(q_i^{b;o,X}(c_d)\) is a probability measure on \(\mathsf K_{i,d}\). The continuum sections and the
finite recognition law are separately declared objects, related only by the hypothesis below.

\subsection{Continuum sections versus finite evaluation}
\label{sec:prob-compatibility}

Each state or model-coordinate law at \((i,d)\) can be obtained in two ways: by sampling the section
at \(c_d\),
\begin{equation}
\operatorname{Samp}_{D,i}\bigl(q_i^{x;o,X}\bigr)=\bigl(q_i^{x;o,X}(c_d)\bigr)_{d\in D_i},
\qquad x\in\{b,m\},
\label{eq:prob-sampling-map}
\end{equation}
or by marginalizing the finite law. The sampling map acts on the response spaces
\(\mathfrak S_i^x\), the marginal maps act on \(\mathcal P(\mathsf Y_D,\mathscr Y_D)\), and neither is
defined through the other.

\propositionheading{Compatibility determines neither side}{prop:prob-compatibility-nonidentifiability}
Suppose the compatibility relation
\begin{equation}
q_i^{b;o,X}(c_d)=q_{i,d}^{b;o,X},
\qquad
q_i^{m;o,X}(c_d)=q_{i,d}^{m;o,X}
\qquad\text{for every }(i,d)\in I_D
\label{eq:prob-compatibility}
\end{equation}
holds. (i) When at least two local agent-block laws are non-Dirac, \eqref{eq:prob-compatibility}
does not determine \(\mathbb Q_{V,o,X}\). (ii) If some support satisfies
\(\mathcal C_i\not\subseteq D\) and the selected fiber of channel \(x\) carries a global chart onto
\(\R^n\) with \(n\geq1\), and the admitted section class is closed under the smooth local
chart perturbations supported off \(D\) used below, then \eqref{eq:prob-compatibility} does not
determine the sections \(q_i^{x;o,X}\). This closure holds for the unrestricted smooth section
class; it need not hold for a declared response subfamily. (iii) If \(\mathcal C_i\subseteq D\) for every
\(i\), which forces every nonempty
support to be zero dimensional, the compatible values determine the sections pointwise.
\status{ESTABLISHED}

\noindent\emph{Proof.}
(i) By \Cref{prop:prob-marginals-do-not-determine-joint} there are distinct laws on
\(\mathsf Y_D\) with the same coordinate marginals, and \eqref{eq:prob-compatibility} constrains the
joint only through those marginals. (ii) Fix \(c^\star\in\mathcal C_i\setminus D\), a neighborhood
\(W\) of \(c^\star\) disjoint from the finite set \(D\) over which \(\mathcal E_x\) is trivialized, the
chart \(\psi\) of the fiber, and a smooth \(\chi:\mathcal C_i\to[0,1]\) supported in \(W\) with
\(\chi(c^\star)=1\); when \(c^\star\) is isolated, \(W=\{c^\star\}\). For \(v\neq0\) replace the
local coordinate \(\beta(c)\) of the section by \(\psi^{-1}(\psi(\beta(c))+\chi(c)v)\) on \(W\). The
result belongs to the admitted class by the local-perturbation hypothesis, differs at
\(c^\star\), and agrees on \(D\), so it satisfies \eqref{eq:prob-compatibility} against the same
marginals. (iii) Equality at every point of
the support is equality of sections. A nonempty open set contained in a finite set consists of
isolated points. \hfill$\square$

\hypothesisheading{Sampling and marginalization are compatible}{hyp:prob-sampling-compatibility}
Where a result transfers a statement between the continuum sections and the finite recognition law,
relation \eqref{eq:prob-compatibility} is adopted. The hypothesis excludes configurations in which an
agent's section and the finite law disagree at a design point, which is possible because the two
objects are declared independently. A result stated entirely about the finite construction does not
use it. Under the respective hypotheses of \Cref{prop:prob-compatibility-nonidentifiability},
adopting it reconstructs neither the joint from its marginals nor a section from finitely many
values; a restricted response family can be identifiable from finitely many evaluations.

\section{What the finite construction does not supply}
\label{sec:prob-continuum}

\openproblemheading{Continuum law theory}{open:prob-continuum-theory}
The finite construction supplies probability spaces and kernels for a declared design \(D\). It
supplies no probability measure on a space of sections over \(\mathcal C\), no unique extension of the
finite marginals to such a space, and no continuum limit, and finite-dimensional consistency alone
does not produce them. Construct a refining family of designs; a projective system of normalized
finite laws indexed by them, so that a Kolmogorov-type extension theorem applies
\citep{Kallenberg2021}; a declared topology and sigma-algebra on a space of sections, with tightness
that places the extension measure on that space; convergence of the evidence bound, the connections,
and the observables that the finite theory computes; control of the gauge choices and of the
reference measures under refinement, neither of which is canonical; and a normalizable continuum
interaction energy.

The last item is where the obstruction is concrete. The reference measures of
\Cref{sec:prob-reference} are finite products, and their sigma-finiteness is inherited from finitely
many sigma-finite factors. The same construction does not extend to infinitely many
non-probability factors, so a continuum theory must declare a measure on the section space directly.
Until that is done a continuum energy has no density to be the exponent of, and every evidence-bound
statement of the main manuscript and of this supplement is a statement about a finite design.

\section{The multi-point generative law}
\label{sec:gen-multi-point-law}

On a design the structural data contain a finite directed acyclic graph \(\Gamma=(V,E)\) whose
incidence is parent closed: if \((i,d)\in I_D\) and \(j\in\operatorname{pa}(i)\), then
\((j,d)\in I_D\). The induced graph \(\Gamma_d\) on \(V_d\) therefore contains every parent of each
participating vertex and is acyclic. At each design point the one-point latent law
\(\mathbb P_{\theta,V,d}^{Y}(dy_d\given X)\) is the directed law of the main manuscript on the agents
\(V_d\), composed from the incidence mechanisms \(G_{\theta,i,d}\) of \eqref{eq:gen-agent-mechanism}
(\Cref{cons:gen-finite-directed-law}).

\hypothesisheading{Design-point product}{hyp:gen-design-product}
Distinct design points are conditionally independent given \(X\):
\begin{equation}
\mathbb P_{\theta,V}^{Y}(dy\given X)
=\bigotimes_{d=1}^{M}\mathbb P_{\theta,V,d}^{Y}(dy_d\given X).
\label{eq:gen-design-product}
\end{equation}
Equivalently, each agent contributes the product mechanism
\begin{equation}
G_{\theta,i,D}:=\bigotimes_{d\in D_i}G_{\theta,i,d}.
\label{eq:gen-design-agent-mechanism}
\end{equation}
The hypothesis excludes residual dependence across design points.

When the component kernels have densities against the declared component reference measures, the
latent density is the product over design points of the one-point densities of
\eqref{eq:gen-density}, \(p_{\theta,V}^Y(y\given X)=\prod_{d=1}^{M}\prod_{i\in V_d}g_{\theta,i,d}\).

\corollaryheading{Factorwise decomposition over a design}{cor:gen-design-factorwise-decomposition}
Assume \Cref{hyp:gen-design-product,hyp:elbo-evidence-domain}, and assume the factorwise log-integrability of
\cref{hyp:elbo-factorwise-integrability} for every design-point factor and every record factor.
Write \(E_{a,o}\) for the negative log density of the record kernel of \(a\in\mathfrak A\) and
\(\partial a\) for its scope. Then
\begin{align*}
\Fenergy_V[\mathbb Q_{V,o,X};X,o]
={}&\sum_{d=1}^{M}\sum_{i\in V_d}
\E_{\mathbb Q_{V,o,X}}\bigl[
-\log p_{\theta,i,d}^{m}\bigl(m_{i,d}\given m_{\operatorname{pa}(i),d},X\bigr)\bigr]\\
&+\sum_{d=1}^{M}\sum_{i\in V_d}
\E_{\mathbb Q_{V,o,X}}\bigl[
-\log p_{\theta,i,d}^{k}\bigl(k_{i,d}\given k_{\operatorname{pa}(i),d},m_{i,d},X\bigr)\bigr]\\
&+\sum_{a\in\mathfrak A}\E_{\mathbb Q_{V,o,X}}\bigl[E_{a,o}(Y_{\partial a})\bigr]
+\E_{\mathbb Q_{V,o,X}}\bigl[\log q_{V,o,X}(Y)\bigr],
\end{align*}
where at a root the parent tuples are empty and the state factor is the model-to-belief bridge. The
generative sectors contain \(2|I_D|+|\mathfrak A|\) factors, which equals \(2MN+|\mathfrak A|\) on a
rectangular design, and the last term is one joint recognition log density, not a sum of marginal
entropies. \status{ESTABLISHED}

\noindent\emph{Proof.}
Under \Cref{hyp:gen-design-product} the complete log density is the sum over design points of the
one-point latent log densities plus the record log densities, and each one-point term is the sum of
its model and state factors. The integrability assumption allows each expectation to be taken
separately, so the argument of \Cref{prop:elbo-factorwise-decomposition} applies term by term, with
the one joint recognition term appearing once on each side. \hfill$\square$

The design-point product carries one copy of each link per design point. For an edge copy
\(e:j\to i\) put \(D_e:=D_i\cap D_j\); the structural data carry copies \(\Theta_{d,e}^x\in G\)
exactly for \(d\in D_e\). Independent frame re-expressions at the incidences,
\(g_{d,i}^x:=g_i^x(c_d)\), act on each copy by the vertex law of the main manuscript,
\(\Theta_{d,e}^{x\prime}=g_{d,i}^x\Theta_{d,e}^x(g_{d,j}^x)^{-1}\) (\Cref{eq:gen-gauge-links}). At one
base point a constraint tying copies to one link is preserved automatically, because all copies
transform by the same pair (\Cref{eq:gen-shared-link-admissibility}); across design points the copies of
one edge transform by different pairs, and preservation becomes a condition.

\definitionheading{Shared-link admissibility and the setwise stabilizer}{def:gen-shared-link-admissibility}
A submodel may constrain the link of a constrained edge copy \(e\) to be one value \(\Theta_e^x\) at
every point of \(D_e\). At a fixed shared-link datum, a rechoice is admissible when for every
constrained \(e\) and all \(d,d'\in D_e\)
\begin{equation}
g_{d,i}^x\Theta_e^x(g_{d,j}^x)^{-1}
=g_{d',i}^x\Theta_e^x(g_{d',j}^x)^{-1}.
\label{eq:gen-design-shared-link-admissibility}
\end{equation}
The admissible family of the datum is the solution set of \eqref{eq:gen-design-shared-link-admissibility}.
The setwise stabilizer is the set of rechoices that map the constraint set, the set of all link
assignments whose constrained links are independent of \(d\in D_e\), into itself.

The admissible family depends on the datum, so admissible rechoices compose as arrows of the action
groupoid of rechoices on link data: \(g\) may follow \(\widetilde g\) when \(g\) is admissible at the
transformed datum \(\widetilde g\cdot\Theta\), not at \(\Theta\). The group that acts on the whole
constrained submodel is the setwise stabilizer, and a larger transformation law requires an
explicitly declared parameter compensation.

\propositionheading{Setwise stabilizer of shared links}{prop:gen-shared-link-stabilizer}
Fix a channel \(x\), let \(Z(G)\) be the center of \(G\), and suppose every constrained link ranges
over all of \(G\). A rechoice \((g_{d,i}^x)\) lies in the setwise stabilizer exactly when, for every
constrained edge copy \(e:j\to i\) and all \(d,d'\in D_e\),
\begin{equation}
(g_{d',i}^x)^{-1}g_{d,i}^x=(g_{d',j}^x)^{-1}g_{d,j}^x\in Z(G).
\label{eq:gen-shared-link-stabilizer}
\end{equation}
Equivalently, on the design points of each constrained edge the rechoice is a context-independent
pair \((\bar g_i^x,\bar g_j^x)\) multiplied by a context-dependent central twist common to both
endpoints, \(g_{d,i}^x=\bar g_i^xz_d\) and \(g_{d,j}^x=\bar g_j^xz_d\) with \(z_d\in Z(G)\). The
context-independent rechoices are the subgroup with trivial twists, and when \(Z(G)\) is trivial the
stabilizer consists exactly of the rechoices that are constant over each \(D_e\) of a constrained
edge. \status{ESTABLISHED}

\noindent\emph{Proof.}
The rechoice maps the constraint set into itself exactly when
\(g_{d,i}^x\Theta(g_{d,j}^x)^{-1}=g_{d',i}^x\Theta(g_{d',j}^x)^{-1}\) for every constrained edge, all
\(d,d'\in D_e\), and every \(\Theta\in G\), that is, when \(\lambda_i\Theta=\Theta\lambda_j\) for all
\(\Theta\), with \(\lambda_i=(g_{d',i}^x)^{-1}g_{d,i}^x\) and
\(\lambda_j=(g_{d',j}^x)^{-1}g_{d,j}^x\). Taking \(\Theta=e_G\) gives \(\lambda_i=\lambda_j\), and
then \(\lambda_i\) commutes with every element of \(G\). Conversely
\(\lambda_i=\lambda_j\in Z(G)\) gives \(\lambda_i\Theta=\Theta\lambda_j\). Fixing \(d'\in D_e\) and
setting \(\bar g_v^x=g_{d',v}^x\) and \(z_d=\lambda_i\) gives the product form. The set is a group,
since for two rechoices with twists \(z_d\) and \(\widetilde z_d\) the product has twist
\(z_d\widetilde z_d\) at both endpoints because central elements commute past
\(\widetilde g_{d',v}^x\), and the inverse has twist \(z_d^{-1}\); so mapping the constraint set into
itself already gives setwise invariance. \hfill$\square$

\propositionheading{Shared-link admissibility need not be closed under composition}{prop:gen-shared-link-nonclosure}
For \(G=\GL(2,\R)\), two design points at which both endpoints participate, one constrained edge,
and \(\Theta_e=I\), the solution set of \eqref{eq:gen-design-shared-link-admissibility} is not closed under
pointwise multiplication. \status{ESTABLISHED}

\noindent\emph{Proof.}
Let \(S=\operatorname{diag}(2,1)\) and \(C=\begin{psmallmatrix}1&1\\0&1\end{psmallmatrix}\). The
rechoice \(g_{1,i}=g_{1,j}=I\), \(g_{2,i}=g_{2,j}=S\) sends the shared link to \(I\) at both design
points, and \(\widetilde g_{d,i}=C\), \(\widetilde g_{d,j}=I\) sends it to \(C\) at both. Their
pointwise product sends it to \(C\) and to \(SCS^{-1}=\begin{psmallmatrix}1&2\\0&1\end{psmallmatrix}\).
\hfill$\square$

The nonclosure is the composition law of the action groupoid: the product acts as
\(g\cdot(\widetilde g\cdot\Theta)\), and \(g\) is admissible at \(I\) but not at \(C\). The first
rechoice is not in the setwise stabilizer because \(S\) is not central. In the Gaussian realization
of Supplement G the links enter the kernels only through their represented copies, so the same
stabilizer acts on a Gaussian submodel with shared links.

\definitionheading{Passive coordinate product}{def:gen-product-action}
At the unconstrained finite-design level, independently declared coordinate re-expressions form the
passive product \(\prod_{(i,d)\in I_D}(G\times G)\) of pairs \((g_{d,i}^b,g_{d,i}^m)\). It records two
independently chosen trivializations of bundles associated to one \(\mathscr P_G\) and is not a
second principal structure group. A context-independent re-expression is an element for which
\(g_{d,i}^x\) is constant as \(d\) ranges over \(D_i\).

\propositionheading{Images in the passive product}{prop:gen-passive-product-images}
(i) If the re-expressions must arise as values \(g_{d,i}^x=g_i^x(c_d)\) of smooth maps
\(g_i^x:\mathcal C_i\to G\), the realized group is the image of this evaluation, and it is the whole
passive product when \(G\) is connected. (ii) The local functions at the design points of the gauge
automorphisms of \(\mathscr P_G\) form a subgroup of \(\prod_{d=1}^{M}\mathcal A_{c_d}\), with
\(\mathcal A_{c_d}\) as in \Cref{prop:geo-two-atlas-active-image}, equal to it when \(G\) is connected
and possibly strictly smaller when \(G\) is disconnected. (iii) For nontrivial \(G\) this active image
is a proper subgroup of the passive product, and a proper subgroup of the product of pointwise
\(h_i\)-twisted diagonals as soon as one design point carries two agents. \status{ESTABLISHED}

\noindent\emph{Proof.}
(i) Given values at the distinct points \(c_d\), the bump-function construction in the proof of
\Cref{prop:geo-fiber-automorphism-realization}, applied to \(G\)-valued maps in disjoint coordinate
neighborhoods, interpolates them, and at isolated points any values are smooth. (ii) The containment
is \Cref{prop:geo-two-atlas-active-image} at each design point; equality for connected \(G\) is
\Cref{prop:geo-fiber-automorphism-realization}, and the \(S_3\)-bundle there with its base point in
\(D\) gives strictness. (iii) Each factor \(\mathcal A_{c_d}\cong G\) is proper in
\(\prod_{i\in V_d}(G\times G)\), and in the product of twisted diagonals when \(|V_d|\geq2\), by
\Cref{prop:geo-two-atlas-active-image}. \hfill$\square$

\openproblemheading{Regular frame-coordinate quotient}{open:gen-gauge-quotient}
Determine the generic stabilizers of the passive \(G\times G\) action generated by independent belief
and model trivializations on the full parameter space, and identify strata on which its orbit space
is a regular quotient manifold. The coordinate action is not the principal structure group, and it
can have nontrivial stabilizers, so subtracting the dimension of the coordinate group from a
parameter count is unjustified without freeness and properness. Nothing in the main manuscript
depends on this quotient.

\section{Scale maps over base maps}
\label{sec:rg-scale-maps-base}

The scale maps of the main manuscript at \(c_*\) are torsor maps fixed by the image of one point
(\Cref{eq:rg-principal-scale-map}). Over a base map their existence is a topological condition. Let
\(\pi_\ell:\mathscr P_\ell\to\mathcal C_\ell\) be principal \(G_\ell\)-bundles over smooth level bases,
\(f_\ell:\mathcal C_\ell\to\mathcal C_{\ell+1}\) smooth, and \(\kappa_\ell:G_\ell\to G_{\ell+1}\) a
Lie-group homomorphism.

\propositionheading{Existence of a principal scale map over a base map}{prop:rg-principal-scale-map-existence}
A map \(\mathcal P_\ell:\mathscr P_\ell\to\mathscr P_{\ell+1}\) covering \(f_\ell\) with
\(\mathcal P_\ell(p\cdot g)=\mathcal P_\ell(p)\cdot\kappa_\ell(g)\) exists if and only if the extended
bundle \(\mathscr P_\ell\times_{\kappa_\ell}G_{\ell+1}\) is isomorphic to the pullback
\(f_\ell^*\mathscr P_{\ell+1}\) as a principal \(G_{\ell+1}\)-bundle over \(\mathcal C_\ell\). The
condition can fail: with the Hopf \(\mathrm U(1)\)-bundle over \(S^2\) as source, the trivial
bundle as target, \(f_\ell=\operatorname{id}\), and \(\kappa_\ell(z)=z^n\) with \(n\neq0\), no such
map exists, while for \(n=0\) the map \(p\mapsto(\pi_\ell(p),1)\) exists for every source bundle.
\status{ESTABLISHED}

\noindent\emph{Proof.}
Given \(\mathcal P_\ell\), the map \([p,g']\mapsto(\pi_\ell(p),\mathcal P_\ell(p)g')\) is well defined,
since \([p\cdot g,\kappa_\ell(g)^{-1}g']\) has the same image, and it is a
\(G_{\ell+1}\)-equivariant map of principal bundles over the identity, hence an isomorphism.
Conversely, an isomorphism \(\Xi\) gives \(\mathcal P_\ell(p)=\operatorname{pr}_2\Xi([p,e])\), and
\([p\cdot g,e]=[p,e]\cdot\kappa_\ell(g)\) gives the equivariance. For the Hopf bundle, the extension
by \(z\mapsto z^n\) is the frame bundle of the \(n\)-th tensor power of the associated line bundle,
whose first Chern number is \(n\) times that of the Hopf bundle, which is nonzero; the trivial bundle
has Chern number zero. For \(n=0\) the extension is trivialized by \([p,g']\mapsto(\pi_\ell(p),g')\).
\hfill$\square$

Given \(\mathcal P_\ell\), a fiber map \(q_{\ell,x}\) descends to a bundle map
\(C_{\ell,x}:\mathcal E_{\ell,x}\to\mathcal E_{\ell+1,x}\) covering \(f_\ell\) exactly under the
intertwining condition \eqref{eq:rg-scale-intertwiner}, by the representative computation that
proves \eqref{eq:rg-associated-scale-map}, which is pointwise in the base. For a multi-step scale
functor the base maps compose as well, \(f_{\ell+1}\circ f_\ell\), and the composition laws of the main
manuscript hold with them.

\propositionheading{Relative frames across scales}{prop:rg-relative-frame-naturality}
Let \(h_\ell\) and \(h_{\ell+1}\) be the relative frames \eqref{eq:rg-relative-frame} of an agent at
levels \(\ell\) and \(\ell+1\), and let comparison functions \(\varsigma_\ell^x\) satisfy
\(\mathcal P_\ell(u_\ell^x(c))=u_{\ell+1}^x(f_\ell c)\cdot\varsigma_\ell^x(c)\). Then
\begin{equation}
\varsigma_\ell^b\kappa_\ell(h_\ell)
=(h_{\ell+1}\circ f_\ell)\varsigma_\ell^m.
\label{eq:rg-relative-frame-naturality}
\end{equation}
\status{ESTABLISHED}

\noindent\emph{Proof.}
\(\mathcal P_\ell(u_\ell^m)=\mathcal P_\ell(u_\ell^bh_\ell)=u_{\ell+1}^b\varsigma_\ell^b\kappa_\ell(h_\ell)
=u_{\ell+1}^m\cdot(h_{\ell+1}\circ f_\ell)^{-1}\varsigma_\ell^b\kappa_\ell(h_\ell)\), and uniqueness in
the torsor against \(\mathcal P_\ell(u_\ell^m)=u_{\ell+1}^m\varsigma_\ell^m\) gives the identity.
\hfill$\square$

The two represented frame channels therefore remain compatible across scales even when their fiber
dimensions or actions differ.

\definitionheading{Connection defects along base curves}{def:rg-connection-defects}
Let each level carry principal connections \(\omega_{\ell,b},\omega_{\ell,m}\) with associated
transports \(\Omega_{\ell,\gamma},\widetilde\Omega_{\ell,\gamma}\), and cross morphisms
\(\Phi_\ell,\widetilde\Phi_\ell\). For a base curve \(\gamma:c\to c'\), the cross-scale connection
defects are the ordered pairs
\begin{align}
\mathbf D^b_{\ell,\gamma}
&=\left(C_{\ell,b,c'}\circ\Omega_{\ell,\gamma},
\Omega_{\ell+1,f_\ell\gamma}\circ C_{\ell,b,c}\right),\nonumber\\
\mathbf D^m_{\ell,\gamma}
&=\left(C_{\ell,m,c'}\circ\widetilde\Omega_{\ell,\gamma},
\widetilde\Omega_{\ell+1,f_\ell\gamma}\circ C_{\ell,m,c}\right),
\label{eq:rg-cross-connection-defects}
\end{align}
and the within-scale nonparallelness pairs are
\begin{align}
\mathbf A^\Phi_{\ell,\gamma}
&=\left(\widetilde\Omega_{\ell,\gamma}\circ\Phi_{\ell,c},
\Phi_{\ell,c'}\circ\Omega_{\ell,\gamma}\right),\nonumber\\
\mathbf A^{\widetilde\Phi}_{\ell,\gamma}
&=\left(\Omega_{\ell,\gamma}\circ\widetilde\Phi_{\ell,c},
\widetilde\Phi_{\ell,c'}\circ\widetilde\Omega_{\ell,\gamma}\right).
\label{eq:rg-within-connection-defects}
\end{align}
A defect vanishes when the two components of its pair are equal. Subtraction of the two components
is reserved for a declared vector or affine realization.

\corollaryheading{Covariance of the connection defects}{cor:rg-connection-defect-covariance}
Under admissible local-frame changes of the belief and model bundles at every level, both components
of every pair in \eqref{eq:rg-cross-connection-defects} and \eqref{eq:rg-within-connection-defects}
acquire the same endpoint coordinate changes, so the equality locus of each pair is gauge invariant,
and in a fiberwise-linear realization their difference obeys the left and right law of a Hom-bundle
section. \status{ESTABLISHED}

\noindent\emph{Proof.}
The transports change by \eqref{eq:geo-pt-gauge-b} and \eqref{eq:geo-pt-gauge-m}, the cross-scale
maps by the endpoint coordinate changes at \(c\) and \(f_\ell(c)\), and the cross morphisms by
\Cref{prop:geo-map-bundle-conversion}; substitution is the computation of
\Cref{prop:rg-defect-covariance}, applied to the connection pairs. \hfill$\square$

\propositionheading{A principal sufficient condition}{prop:rg-principal-connection-naturality}
If
\begin{equation}
\mathcal P_\ell^*\omega_{\ell+1,x}
=\mathrm d\kappa_\ell\circ\omega_{\ell,x},
\label{eq:rg-principal-connection-naturality}
\end{equation}
then \(\mathbf D^x_{\ell,\gamma}\) vanishes for every base curve. If \(\rho_x\) is not faithful,
vanishing of the represented defect does not imply \eqref{eq:rg-principal-connection-naturality}.
\status{ESTABLISHED}

\noindent\emph{Proof.}
Under \eqref{eq:rg-principal-connection-naturality},
\(\omega_{\ell+1,x}(T\mathcal P_\ell Y)=\mathrm d\kappa_\ell(\omega_{\ell,x}(Y))=0\) for horizontal
\(Y\), so \(\mathcal P_\ell\) carries horizontal lifts of \(\gamma\) to horizontal lifts of
\(f_\ell\gamma\), and the associated maps intertwine the transports. For the trivial representation
every represented transport is the identity, so the represented defect vanishes for all connections,
while \eqref{eq:rg-principal-connection-naturality} can fail. \hfill$\square$

The infinitesimal failure of \eqref{eq:rg-principal-connection-naturality} is the scale-connection
defect form \eqref{eq:pb-scale-connection-defect}; the represented horizontal defect vanishes along a
coarse section exactly under the isotropy criterion of \Cref{thm:pb-isotropy-criterion}, and the
defects of composable steps obey the ordered law \eqref{eq:rg-scale-connection-defect-cocycle}.

An associated-bundle isomorphism \(\Xi_\ell:\mathcal E_{\ell,b}\to\mathcal E_{\ell,m}\) may also be
declared. Its existence is additional global data: equal fiber type is necessary but not sufficient,
as the trivial and M\"obius line bundles of \Cref{prop:geo-cech-class} show, and \(\Xi_\ell\) is
unavailable when the associated bundles are not isomorphic. Its parallel and scale-preservation
conditions are \(\widetilde\Omega_{\ell,\gamma}\Xi_{\ell,c}=\Xi_{\ell,c'}\Omega_{\ell,\gamma}\) and
\(C_{\ell,m}\Xi_\ell=\Xi_{\ell+1}C_{\ell,b}\). No such isomorphism is needed to define
\(\Phi_\ell\) and \(\widetilde\Phi_\ell\).

\section{Variable contexts and ordered Kron frames}
\label{sec:cg-kron-variable-context}

Supplement G classifies the Gaussian coefficient families that pass the fixed-anchor pair test at one
anchor (\Cref{thm:cg-independent-kron-classification}). When the anchor varies over contexts, the
classification holds locally and its ordering can fail globally.

\theoremheading{Local frames and a global ordering obstruction}{thm:cg-kron-local-context-frame}
On a smooth context domain \(\mathcal U\subseteq\mathcal C\), let \(Q_0(c)\succ0\) vary smoothly, let
\(\mathscr W_c\) be an inclusion-maximal family for the fixed-anchor pair test at \(Q_0(c)\), and
suppose that the whitened spans
\begin{equation}
\mathscr D_c
=\operatorname{span}_{\R}
\{Q_0(c)^{-1/2}AQ_0(c)^{-1/2}:A\in\mathscr W_c\}
\label{eq:cg-kron-context-algebra}
\end{equation}
form a smooth rank-\(K\) vector subbundle of \(\mathcal U\times\Sym^K\); no pair using coefficients
from different contexts is admitted. Then every \(c_0\in\mathcal U\) has a neighborhood
\(\mathcal U_0\) and a smooth \(O:\mathcal U_0\to\mathrm O(K)\) with
\begin{equation}
\begin{aligned}
\mathscr D_c&=O(c)\operatorname{Diag}(K)O(c)^\top,\\
\mathscr W_c&=\mathsf Q(c)\operatorname{Diag}_+(K)\mathsf Q(c)^\top,\\
\mathsf Q(c)&=Q_0(c)^{1/2}O(c).
\end{aligned}
\label{eq:cg-kron-local-context-cone}
\end{equation}
A global continuous ordered diagonal frame need not exist, even for a smooth maximal family over the
circle with \(Q_0=I_2\). Smooth global coefficients nevertheless give smooth Schur complements and
log determinants wherever the precision blocks remain positive definite, and if every cross-context
pair were tested against one common fixed anchor, the fixed-anchor classification would apply to the
union and force one common frame. \status{ESTABLISHED}

\noindent\emph{Proof.}
The implicit-function theorem applied on symmetric matrices to \(S^2-Q_0=0\) at the positive square
root has derivative \(X\mapsto SX+XS\), which multiplies entry \(ij\) by \(s_i+s_j>0\) in an
eigenbasis of \(S\); so \(Q_0^{1/2}\) and its inverse are smooth, and uniqueness of the positive
square root identifies the local solutions. By \Cref{thm:cg-independent-kron-classification},
\(\mathscr D_{c_0}\) contains a matrix with \(K\) distinct eigenvalues, and a smooth local frame of
the subbundle extends it to a smooth section \(\mathsf E(c)\in\mathscr D_c\). For
\(\det(\lambda I-\mathsf E(c))=0\), the derivative in \(\lambda\) at each initial eigenvalue is the
product of its nonzero spectral gaps, so the implicit-function theorem supplies \(K\) distinct smooth
eigenvalues \(\lambda_k(c)\) on a smaller neighborhood, with spectral projectors
\begin{equation}
\mathsf P_k(c)=\prod_{l\ne k}
\frac{\mathsf E(c)-\lambda_l(c)I}{\lambda_k(c)-\lambda_l(c)}.
\label{eq:cg-kron-smooth-projectors}
\end{equation}
In an eigenbasis this polynomial is one on the \(k\)-th eigenspace and zero on the others, so these
are smooth orthogonal rank-one projectors. For a unit \(v_k\in\operatorname{im}\mathsf P_k(c_0)\),
\(e_k(c)=\mathsf P_k(c)v_k/\lVert\mathsf P_k(c)v_k\rVert\) is smooth after shrinking the
neighborhood, and its columns form \(O(c)\). Every member of \(\mathscr D_c\) commutes with
\(\mathsf E(c)\) and preserves its eigenlines, and both spaces have dimension \(K\), so the diagonal
containment is equality; undoing the whitening gives \eqref{eq:cg-kron-local-context-cone}.

For the obstruction write \(t\in S^1=\R/(2\pi\mathbb Z)\), let \(\operatorname{Rot}(\phi)\) be the
planar rotation by \(\phi\), and set
\begin{equation}
\mathscr W_t=
\operatorname{Rot}(t/4)\operatorname{Diag}_+(2)
\operatorname{Rot}(t/4)^\top,\qquad
\mathsf E_t=
\begin{pmatrix}
\cos(t/2)&\sin(t/2)\\
\sin(t/2)&-\cos(t/2)
\end{pmatrix}.
\label{eq:cg-kron-circle-cone}
\end{equation}
Increasing \(t\) by \(2\pi\) rotates by another \(\pi/2\), which permutes the coordinate axes and
preserves the cone. The algebra is \(\operatorname{span}\{I_2,\mathsf E_t\}\), with trace-inner-product
projection
\begin{equation}
X\longmapsto
\frac12\Tr(X)I_2+\frac12\Tr(\mathsf E_tX)\mathsf E_t,
\label{eq:cg-kron-circle-algebra-projector}
\end{equation}
which is smooth and periodic because \(\mathsf E_{t+2\pi}=-\mathsf E_t\), and has constant rank two,
so it defines the required subbundle. The two trace-one extreme projectors
\((I_2\pm\mathsf E_t)/2\) interchange after a full circuit. A continuous ordered first projector
\(\mathsf P(t)\) has \(\Tr(\mathsf E_t\mathsf P(t))\in\{1,-1\}\), a continuous function on
\([0,2\pi]\), so its sign is constant, and its endpoint is the opposite projector from its initial
value, contradicting periodicity; an ordered orthogonal diagonal frame would supply such a projector.
Matrix inversion is the adjugate divided by a nonzero determinant and \(\log\det\) is smooth on
positive-definite matrices, so no diagonal frame is needed for Schur complements or normalizers. The
last statement is \Cref{thm:cg-independent-kron-classification} applied to the union of the
families. \hfill$\square$
